%%%%%%%%%%%%%%%%%%%%%%%%%%%%%%%%%%%%%%%%%%%%%%%%%%%%%%%%%%%%%%%%%%%%%%%%%%%%%%%
\subsection{Awareness as a Unifying Framework for Cognitive Biases}
\label{subsec:awareness-biases}
%%%%%%%%%%%%%%%%%%%%%%%%%%%%%%%%%%%%%%%%%%%%%%%%%%%%%%%%%%%%%%%%%%%%%%%%%%%%%%%

A remarkable property of the Awareness function is its capacity to explain
and systematically organize a wide range of cognitive biases documented in
behavioral economics and psychology. Rather than treating each bias as an
isolated phenomenon, the Awareness framework reveals biases as \emph{systematic
asymmetries} in the six-component awareness structure.

\subsubsection{The Core Insight: Biases as Awareness Asymmetries}

\begin{tcolorbox}[colback=blue!5!white, colframe=blue!75!black, title=Central Thesis]
\textbf{Cognitive biases are not random errors. They are systematic consequences
of asymmetries in the Awareness structure:}
\begin{itemize}
\item Asymmetries between system levels (SYS 0, SYS 1, SYS 2)
\item Asymmetries between self and community awareness
\item Asymmetries between utility categories (INU, KNU, IDN)
\item Asymmetries created by awareness determinants
\end{itemize}
\end{tcolorbox}

This unifying perspective has three implications:
\begin{enumerate}
\item \textbf{Explanatory:} Each bias can be traced to specific awareness asymmetries
\item \textbf{Predictive:} We can predict when biases will occur based on awareness patterns
\item \textbf{Interventional:} Bias correction becomes awareness correction
\end{enumerate}

%%%%%%%%%%%%%%%%%%%%%%%%%%%%%%%%%%%%%%%%%%%%%%%%%%%%%%%%%%%%%%%%%%%%%%%%%%%%%%%
\subsubsection{Category 1: System-Level Biases (SYS 0 vs. SYS 1 vs. SYS 2)}
%%%%%%%%%%%%%%%%%%%%%%%%%%%%%%%%%%%%%%%%%%%%%%%%%%%%%%%%%%%%%%%%%%%%%%%%%%%%%%%

These biases arise from asymmetries between the three processing systems.

\paragraph{The General Pattern.}
\begin{equation}
\text{System-Level Bias} = f(A_{SYS0} - A_{SYS2}) \quad \text{or} \quad f(A_{SYS1} - A_{SYS2})
\end{equation}

When faster systems (SYS 0, SYS 1) dominate slower systems (SYS 2), 
systematic biases emerge.

\begin{table}[htbp]
\centering
\caption{System-Level Biases Explained by Awareness}
\label{tab:system-level-biases}
\renewcommand{\arraystretch}{1.4}
\begin{tabular}{@{}lp{4cm}p{6cm}@{}}
\toprule
\textbf{Bias} & \textbf{Traditional Definition} & \textbf{Awareness Explanation} \\
\midrule
\textbf{Present Bias} & 
Overweighting immediate outcomes & 
$A_{SYS0}(t=0) = 1$ automatic; $A_{SYS2}(t>0)$ requires effort and is typically low \\
\midrule
\textbf{Hot-Cold Empathy Gap} & 
Underestimating visceral influences when ``cold'' & 
When in cold state: $A_{SYS0} \approx 0$, only $A_{SYS2}$ active; when hot: $A_{SYS0} = 1$ dominates \\
\midrule
\textbf{Affect Heuristic} & 
Judgments driven by current feelings & 
$A_{SYS1}(emotion)$ high; $A_{SYS2}(analysis)$ suppressed by Hyper-Dissonance Rule \\
\midrule
\textbf{Hyperbolic Discounting} & 
Inconsistent time preferences & 
$A_{SYS0}$ for immediate rewards; $A_{SYS2}$ for delayed rewards (effortful, inconsistent) \\
\midrule
\textbf{Temptation/Akrasia} & 
Acting against better judgment & 
$A_{SYS0}(temptation) = 1$ instant; $A_{SYS2}(consequences)$ requires activation \\
\midrule
\textbf{Projection Bias} & 
Projecting current state to future & 
Current $A_{SYS0/1}$ dominates; future $A_{SYS2}$ underactivated \\
\midrule
\textbf{Immediacy Effect} & 
Disproportionate weight on immediate & 
$A_{instant}$ automatic (SYS 0); $A_{delayed}$ requires SYS 2 reflection \\
\bottomrule
\end{tabular}
\end{table}

\paragraph{The Hyper-Dissonance Mechanism.}
Recall from Section~\ref{subsec:system-levels}:
\begin{equation}
A_{SYS0}(t^*) = 1 \quad \Rightarrow \quad A_{SYS2}(t^*) < 1
\end{equation}

This constraint explains why visceral states reliably produce biased decisions:
when SYS 0 is maximally activated, reflective awareness is \emph{structurally}
suppressed.

\paragraph{Example: Present Bias in Saving.}
\begin{itemize}
\item $A_{SYS0}$: Immediate consumption pleasure = 1 (automatic)
\item $A_{SYS2}$: Future retirement security = requires effortful activation
\item Result: Systematic undersaving despite knowing better
\end{itemize}

%%%%%%%%%%%%%%%%%%%%%%%%%%%%%%%%%%%%%%%%%%%%%%%%%%%%%%%%%%%%%%%%%%%%%%%%%%%%%%%
\subsubsection{Category 2: Self-Community Biases}
%%%%%%%%%%%%%%%%%%%%%%%%%%%%%%%%%%%%%%%%%%%%%%%%%%%%%%%%%%%%%%%%%%%%%%%%%%%%%%%

These biases arise from the awareness gradient between self and community:
\begin{equation}
A^{self} >> A^{comm}
\end{equation}

\paragraph{The General Pattern.}
\begin{equation}
\text{Self-Community Bias} = f(A^{self}_X - A^{comm}_X)
\end{equation}

Because awareness of effects on others requires deliberate perspective-taking
(especially for $A^{comm}_{KNU}$ and $A^{comm}_{IDN}$), systematic biases emerge.

\begin{table}[htbp]
\centering
\caption{Self-Community Biases Explained by Awareness}
\label{tab:self-community-biases}
\renewcommand{\arraystretch}{1.4}
\begin{tabular}{@{}lp{4cm}p{6cm}@{}}
\toprule
\textbf{Bias} & \textbf{Traditional Definition} & \textbf{Awareness Explanation} \\
\midrule
\textbf{Egocentric Bias} & 
Overweighting own perspective & 
$A^{self}$ high by default; $A^{comm}$ requires activation \\
\midrule
\textbf{Bystander Effect} & 
Not helping when others present & 
$A^{comm}_{INU}$ (victim's welfare) diluted across observers; $A^{self}_{INU}$ (cost of helping) remains high \\
\midrule
\textbf{Diffusion of Responsibility} & 
Reduced sense of personal obligation & 
$A^{self}_{KNU}$ (my responsibility to group) decreases with group size \\
\midrule
\textbf{False Consensus Effect} & 
Assuming others share our views & 
$A^{comm}_{IDN}$ (others' actual identity/views) very low; project $A^{self}_{IDN}$ \\
\midrule
\textbf{Out-Group Homogeneity} & 
Seeing out-groups as uniform & 
$A^{comm}_{IDN}$ for out-group $\approx 0$; no differentiation of their identities \\
\midrule
\textbf{Spotlight Effect} & 
Overestimating others' attention to us & 
High $A^{self}_{IDN}$; assume others have equally high $A^{comm}$ for us \\
\midrule
\textbf{Empathy Gap (Interpersonal)} & 
Failing to understand others' states & 
$A^{comm}_{INU}$ and $A^{comm}_{IDN}$ require deliberate empathy; often not activated \\
\midrule
\textbf{Curse of Knowledge} & 
Assuming others know what we know & 
$A^{self}$ (our knowledge) high; $A^{comm}$ (their knowledge state) not activated \\
\bottomrule
\end{tabular}
\end{table}

\paragraph{The Gradient Mechanism.}
The awareness gradient (Section~\ref{subsec:two-level-structure}) creates 
a \emph{structural} bias toward self-focused decisions:

\begin{equation}
A^{self}_{INU} > A^{self}_{KNU} > A^{self}_{IDN} > A^{comm}_{INU} > A^{comm}_{KNU} > A^{comm}_{IDN}
\end{equation}

Each step down the gradient requires more cognitive effort, making 
self-other asymmetries predictable and systematic.

\paragraph{Example: Bystander Effect.}
\begin{itemize}
\item $A^{self}_{INU}$: Cost of helping (effort, risk) = HIGH
\item $A^{comm}_{INU}$: Victim's suffering = MEDIUM (activated by seeing)
\item With multiple bystanders: $A^{comm}_{INU}$ per person decreases (diffusion)
\item Result: No one helps despite all recognizing the need
\end{itemize}

%%%%%%%%%%%%%%%%%%%%%%%%%%%%%%%%%%%%%%%%%%%%%%%%%%%%%%%%%%%%%%%%%%%%%%%%%%%%%%%
\subsubsection{Category 3: INU-KNU-IDN Biases}
%%%%%%%%%%%%%%%%%%%%%%%%%%%%%%%%%%%%%%%%%%%%%%%%%%%%%%%%%%%%%%%%%%%%%%%%%%%%%%%

These biases arise from asymmetries between utility categories, especially
the critical fact that \textbf{KNU is never instant} (AWX-A37).

\paragraph{The General Pattern.}
\begin{equation}
\text{Category Bias} = f(A_{INU} - A_{KNU}) \quad \text{or} \quad f(A_{IDN} - A_{KNU})
\end{equation}

Because individual utility can be instant but collective utility cannot,
systematic biases toward individual over collective outcomes emerge.

\begin{table}[htbp]
\centering
\caption{INU-KNU-IDN Biases Explained by Awareness}
\label{tab:category-biases}
\renewcommand{\arraystretch}{1.4}
\begin{tabular}{@{}lp{4cm}p{6cm}@{}}
\toprule
\textbf{Bias/Phenomenon} & \textbf{Traditional Definition} & \textbf{Awareness Explanation} \\
\midrule
\textbf{Tragedy of the Commons} & 
Overusing shared resources & 
$A_{INU}$ instant (my consumption); $A_{KNU}$ never instant (commons depletion requires SYS 2) \\
\midrule
\textbf{Free Rider Problem} & 
Benefiting without contributing & 
$A^{self}_{INU}$ (benefit) high; $A^{self}_{KNU}$ (my contribution to collective) low \\
\midrule
\textbf{Social Dilemmas} & 
Individual vs. collective rationality & 
$A_{INU}$ > $A_{KNU}$ structurally; KNU always requires awareness activation \\
\midrule
\textbf{NIMBY} & 
Not In My Back Yard & 
$A^{self}_{INU}$ (local costs) instant; $A^{comm}_{KNU}$ (societal benefits) requires SYS 2 \\
\midrule
\textbf{Identity-Protective Cognition} & 
Rejecting facts threatening identity & 
$A^{self}_{IDN}$ (identity protection) dominates $A^{self}_{INU}$ (factual accuracy) \\
\midrule
\textbf{Motivated Reasoning} & 
Reasoning toward desired conclusions & 
$A^{self}_{IDN}$ or $A^{self}_{INU}$ activates first; subsequent $A_{SYS2}$ is biased \\
\midrule
\textbf{System Justification} & 
Defending existing systems & 
$A^{self}_{IDN}$ (identity tied to system) blocks $A^{comm}_{KNU}$ (system's collective costs) \\
\midrule
\textbf{In-Group Favoritism} & 
Preferring in-group members & 
$A^{self}_{IDN}$ (group identity) activates $A^{self}_{KNU}$ for in-group only \\
\bottomrule
\end{tabular}
\end{table}

\paragraph{The KNU-Never-Instant Mechanism.}
This is perhaps the most consequential asymmetry:

\begin{equation}
U^{INU}_{instant} \Rightarrow A_{INU} = 1 \quad \text{(automatic)}
\end{equation}
\begin{equation}
U^{KNU}_{any} \Rightarrow A_{KNU} < 1 \quad \text{(requires deliberate activation)}
\end{equation}

This structural asymmetry explains why:
\begin{itemize}
\item Environmental behavior is systematically underproduced
\item Public goods are systematically underprovided
\item Commons are systematically overexploited
\item Collective action problems are so persistent
\end{itemize}

\paragraph{Example: Climate Inaction.}
\begin{itemize}
\item $A^{self}_{INU}$: Cost of green behavior (money, convenience) = HIGH (instant)
\item $A^{comm}_{KNU}$: Climate benefit to global community = VERY LOW (requires SYS 2, abstract, distant)
\item Result: Systematic underinvestment in climate action
\end{itemize}

%%%%%%%%%%%%%%%%%%%%%%%%%%%%%%%%%%%%%%%%%%%%%%%%%%%%%%%%%%%%%%%%%%%%%%%%%%%%%%%
\subsubsection{Category 4: Determinant-Based Biases}
%%%%%%%%%%%%%%%%%%%%%%%%%%%%%%%%%%%%%%%%%%%%%%%%%%%%%%%%%%%%%%%%%%%%%%%%%%%%%%%

These biases arise from the 20 determinants (AWX-A11 to AWX-A30) that
shape awareness activation.

\paragraph{Cognitive Determinant Biases.}

\begin{table}[htbp]
\centering
\caption{Cognitive Determinant Biases}
\label{tab:cognitive-biases}
\renewcommand{\arraystretch}{1.4}
\begin{tabular}{@{}llp{5.5cm}@{}}
\toprule
\textbf{Bias} & \textbf{Determinant} & \textbf{Mechanism} \\
\midrule
\textbf{Availability Heuristic} & AWX-A14 & $A \propto Availability$; easily recalled events get higher awareness \\
\textbf{Salience Bias} & AWX-A9 & $A = f(Salience)$; vivid stimuli dominate awareness \\
\textbf{Anchoring} & AWX-A9, A27 & First information primes $A$; subsequent adjustment insufficient \\
\textbf{Framing Effects} & AWX-A47 & Different frames activate different $A$ components \\
\textbf{Base Rate Neglect} & AWX-A14 & Vivid cases ($A$ high) dominate statistical bases ($A$ low) \\
\textbf{Conjunction Fallacy} & AWX-A14 & Detailed scenarios more ``available'' than abstract probabilities \\
\bottomrule
\end{tabular}
\end{table}

\paragraph{Emotional Determinant Biases.}

\begin{table}[htbp]
\centering
\caption{Emotional Determinant Biases}
\label{tab:emotional-biases}
\renewcommand{\arraystretch}{1.4}
\begin{tabular}{@{}llp{5.5cm}@{}}
\toprule
\textbf{Bias} & \textbf{Determinant} & \textbf{Mechanism} \\
\midrule
\textbf{Negativity Bias} & AWX-A15, A16 & Fear (AWX-A16) amplifies $A$ for negative more than hope (AWX-A17) for positive \\
\textbf{Loss Aversion} & AWX-A16 & $A_{Pain}$ (fear of loss) > $A_{Gain}$ (hope of gain) \\
\textbf{Mood-Congruent Recall} & AWX-A15 & Current emotion activates mood-congruent $A$ \\
\textbf{Risk Aversion (Gains)} & AWX-A16 & Fear of losing gain activates $A$ for certain option \\
\textbf{Risk Seeking (Losses)} & AWX-A17 & Hope of avoiding loss activates $A$ for risky option \\
\bottomrule
\end{tabular}
\end{table}

\paragraph{Social Determinant Biases.}

\begin{table}[htbp]
\centering
\caption{Social Determinant Biases}
\label{tab:social-biases}
\renewcommand{\arraystretch}{1.4}
\begin{tabular}{@{}llp{5.5cm}@{}}
\toprule
\textbf{Bias} & \textbf{Determinant} & \textbf{Mechanism} \\
\midrule
\textbf{Conformity Bias} & AWX-A19 & $A \propto Norm_{pressure}$; social norms dominate private judgment \\
\textbf{Bandwagon Effect} & AWX-A20, A21 & $A \propto Peer_{visibility}$, $Network_{exposure}$ \\
\textbf{Authority Bias} & AWX-A22 & $A \propto Institutional_{trust}$; authority activates $A$ \\
\textbf{Social Proof} & AWX-A20 & Others' behavior increases $A$ for same behavior \\
\textbf{Herd Behavior} & AWX-A21 & Network effects amplify $A$ cascades \\
\bottomrule
\end{tabular}
\end{table}

\paragraph{Information and Complexity Biases.}

\begin{table}[htbp]
\centering
\caption{Complexity and Uncertainty Biases}
\label{tab:complexity-biases}
\renewcommand{\arraystretch}{1.4}
\begin{tabular}{@{}llp{5.5cm}@{}}
\toprule
\textbf{Bias} & \textbf{Axiom} & \textbf{Mechanism} \\
\midrule
\textbf{Overconfidence} & AWX-A47a & Failure Awareness Gap: failures invisible → overestimate success \\
\textbf{Survivorship Bias} & AWX-A47a & $A_{success} >> A_{failure}$ due to visibility asymmetry \\
\textbf{Hindsight Bias} & AWX-A14 & Outcomes now ``available''; pre-outcome uncertainty not salient \\
\textbf{Planning Fallacy} & AWX-A51 & Epistemic uncertainty reduces $A$ for obstacles \\
\textbf{Illusion of Control} & AWX-A47a & $A$ for own actions high; $A$ for randomness low \\
\textbf{Ambiguity Aversion} & AWX-A52 & Accumulated uncertainty depresses $A$; prefer known risks \\
\bottomrule
\end{tabular}
\end{table}

%%%%%%%%%%%%%%%%%%%%%%%%%%%%%%%%%%%%%%%%%%%%%%%%%%%%%%%%%%%%%%%%%%%%%%%%%%%%%%%
\subsubsection{Category 5: Temporal and Dynamic Biases}
%%%%%%%%%%%%%%%%%%%%%%%%%%%%%%%%%%%%%%%%%%%%%%%%%%%%%%%%%%%%%%%%%%%%%%%%%%%%%%%

These biases arise from the temporal properties of awareness (AWX-A56 to AWX-A65).

\begin{table}[htbp]
\centering
\caption{Temporal Biases Explained by Awareness}
\label{tab:temporal-biases}
\renewcommand{\arraystretch}{1.4}
\begin{tabular}{@{}llp{5.5cm}@{}}
\toprule
\textbf{Bias} & \textbf{Axiom} & \textbf{Mechanism} \\
\midrule
\textbf{Recency Bias} & AWX-A58 & $A$ decays; recent events have higher residual $A$ \\
\textbf{Primacy Effect} & AWX-A27 & First information primes $A$; anchoring \\
\textbf{Peak-End Rule} & AWX-A63 & Shock moments ($A$ spikes) and end ($A$ at recall) dominate \\
\textbf{Duration Neglect} & AWX-A58 & $A$ integrates peaks, not duration \\
\textbf{Status Quo Bias} & AWX-A10 & Habitualized patterns require no $A$; change requires $A$ activation \\
\textbf{Mere Exposure Effect} & AWX-A59, A60 & Repeated exposure increases $A$ (familiarity) without effortful processing \\
\textbf{Normalcy Bias} & AWX-A60 & Habitual weakening: threats become ``normal,'' $A$ drops \\
\bottomrule
\end{tabular}
\end{table}

\paragraph{The Habituation Mechanism.}
AWX-A10 (Habituation) creates a particularly important class of biases:

\begin{equation}
\text{Habituation} \uparrow \quad \Rightarrow \quad A_X(t^*) \rightarrow 0
\end{equation}

When behaviors become automatic:
\begin{itemize}
\item No awareness of their effects is required
\item No awareness of alternatives is activated
\item Status quo is perpetuated without evaluation
\end{itemize}

This explains status quo bias, default effects, and resistance to change.

%%%%%%%%%%%%%%%%%%%%%%%%%%%%%%%%%%%%%%%%%%%%%%%%%%%%%%%%%%%%%%%%%%%%%%%%%%%%%%%
\subsubsection{Category 6: Conflict and Defensive Biases}
%%%%%%%%%%%%%%%%%%%%%%%%%%%%%%%%%%%%%%%%%%%%%%%%%%%%%%%%%%%%%%%%%%%%%%%%%%%%%%%

These biases arise from the conflict and dissonance mechanisms (AWX-A66 to AWX-A70).

\begin{table}[htbp]
\centering
\caption{Defensive Biases Explained by Awareness}
\label{tab:defensive-biases}
\renewcommand{\arraystretch}{1.4}
\begin{tabular}{@{}llp{5.5cm}@{}}
\toprule
\textbf{Bias} & \textbf{Axiom} & \textbf{Mechanism} \\
\midrule
\textbf{Confirmation Bias} & AWX-A66 & Suppress $A$ for disconfirming information (unpleasant) \\
\textbf{Denial} & AWX-A66 & Complete suppression of $A$ for threatening information \\
\textbf{Cognitive Dissonance Reduction} & AWX-A67 & Split $A_{cognitive}$ and $A_{emotional}$ to reduce discomfort \\
\textbf{Rationalization} & AWX-A66, A67 & Post-hoc $A_{SYS2}$ justification after $A_{SYS0/1}$ decision \\
\textbf{Self-Serving Bias} & AWX-A66 & Suppress $A$ for own responsibility in failures \\
\textbf{Defensive Attribution} & AWX-A66 & $A$ for internal causes when outcome is positive; external when negative \\
\textbf{Belief Perseverance} & AWX-A66, A69 & Once $A$ pattern established, contradictory info suppressed \\
\bottomrule
\end{tabular}
\end{table}

\paragraph{The Suppression Mechanism.}
AWX-A66 describes how unpleasant awareness is systematically suppressed:

\begin{equation}
\text{If } A_X(t^*) \rightarrow \text{strong negative emotion} \quad \Rightarrow \quad A_X(t^*) \downarrow
\end{equation}

This creates a \emph{motivated} asymmetry in awareness:
\begin{itemize}
\item $A$ for ego-consistent information: HIGH
\item $A$ for ego-threatening information: SUPPRESSED
\end{itemize}

\paragraph{The Mirror Gap Mechanism.}
AWX-A70a explains defensive reactions to self-image threats:

\begin{equation}
\Delta I_t \geq 0.6 \quad \Rightarrow \quad \frac{\partial A_X}{\partial \Delta I_t} < 0 \quad \text{(defensive)}
\end{equation}

When the gap between self-image and social feedback becomes large,
awareness \emph{decreases} rather than increases---a protective mechanism
that perpetuates biased self-perception.

%%%%%%%%%%%%%%%%%%%%%%%%%%%%%%%%%%%%%%%%%%%%%%%%%%%%%%%%%%%%%%%%%%%%%%%%%%%%%%%
\subsubsection{The Complete Bias-Awareness Map}
%%%%%%%%%%%%%%%%%%%%%%%%%%%%%%%%%%%%%%%%%%%%%%%%%%%%%%%%%%%%%%%%%%%%%%%%%%%%%%%

\begin{table}[htbp]
\centering
\caption{Summary: 40+ Biases Mapped to Awareness Mechanisms}
\label{tab:complete-bias-map}
\small
\renewcommand{\arraystretch}{1.2}
\begin{tabular}{@{}llll@{}}
\toprule
\textbf{Bias Category} & \textbf{Count} & \textbf{Primary Mechanism} & \textbf{Key Axioms} \\
\midrule
System-Level (SYS 0/1/2) & 7 & SYS 0/1 dominates SYS 2 & Hyper-Dissonance \\
Self-Community & 8 & $A^{self} >> A^{comm}$ & Gradient \\
INU-KNU-IDN & 8 & KNU never instant & AWX-A37 \\
Cognitive Determinants & 6 & Availability, Salience & AWX-A14, A9 \\
Emotional Determinants & 5 & Fear, Hope asymmetry & AWX-A15-17 \\
Social Determinants & 5 & Norm pressure, Network & AWX-A19-21 \\
Complexity/Uncertainty & 6 & Failure Gap, Uncertainty & AWX-A47a, A51 \\
Temporal/Dynamic & 7 & Decay, Habituation & AWX-A58, A60 \\
Defensive/Conflict & 7 & Suppression, Mirror Gap & AWX-A66, A70a \\
\midrule
\textbf{Total} & \textbf{59} & & \\
\bottomrule
\end{tabular}
\end{table}

%%%%%%%%%%%%%%%%%%%%%%%%%%%%%%%%%%%%%%%%%%%%%%%%%%%%%%%%%%%%%%%%%%%%%%%%%%%%%%%
\subsubsection{Implications for Debiasing}
%%%%%%%%%%%%%%%%%%%%%%%%%%%%%%%%%%%%%%%%%%%%%%%%%%%%%%%%%%%%%%%%%%%%%%%%%%%%%%%

If biases are awareness asymmetries, then debiasing is awareness correction.

\paragraph{The Debiasing Principle.}
\begin{equation}
\text{Debias} = \text{Correct } A_{\text{underactivated}} \text{ relative to } A_{\text{overactivated}}
\end{equation}

\begin{table}[htbp]
\centering
\caption{Debiasing Strategies Through Awareness Correction}
\label{tab:debiasing}
\renewcommand{\arraystretch}{1.4}
\begin{tabular}{@{}lll@{}}
\toprule
\textbf{Bias Type} & \textbf{Asymmetry} & \textbf{Debiasing Strategy} \\
\midrule
System-Level & SYS 0 > SYS 2 & Create ``cooling off'' periods; activate SYS 2 before SYS 0 \\
Self-Community & $A^{self} > A^{comm}$ & Perspective-taking exercises; identifiable victims \\
INU-KNU & $A_{INU} > A_{KNU}$ & Make collective effects visible; social proof \\
Availability & Available > Base rate & Provide statistical information; slow down \\
Emotional & Fear > Hope & Balance emotional framing \\
Social & Norm > Private & Create safe spaces for dissent \\
Defensive & Confirming > Disconfirming & Pre-commit to considering alternatives \\
\bottomrule
\end{tabular}
\end{table}

\paragraph{The Awareness-Based Debiasing Protocol.}
\begin{enumerate}
\item \textbf{Identify the bias:} What systematic error is occurring?
\item \textbf{Map to awareness:} Which $A$ component is over/underactivated?
\item \textbf{Diagnose the mechanism:} System-level? Self-Community? Determinant?
\item \textbf{Design intervention:} Target the specific awareness asymmetry
\item \textbf{Verify:} Did the asymmetry decrease?
\end{enumerate}

\paragraph{Example: Debiasing Overconfidence.}
\begin{enumerate}
\item \textbf{Bias:} Overconfidence in predictions
\item \textbf{Awareness map:} $A_{success} >> A_{failure}$ (Failure Gap)
\item \textbf{Mechanism:} AWX-A47a---failures invisible due to visibility cost, power distance
\item \textbf{Intervention:} 
\begin{itemize}
\item Increase psychological safety (reduce visibility cost)
\item Systematically surface failures (reduce $\Delta A_F$)
\item ``Pre-mortem'' exercises (activate $A$ for failure scenarios)
\end{itemize}
\item \textbf{Verify:} Calibration improves; prediction intervals widen appropriately
\end{enumerate}

%%%%%%%%%%%%%%%%%%%%%%%%%%%%%%%%%%%%%%%%%%%%%%%%%%%%%%%%%%%%%%%%%%%%%%%%%%%%%%%
\subsubsection{Theoretical Integration}
%%%%%%%%%%%%%%%%%%%%%%%%%%%%%%%%%%%%%%%%%%%%%%%%%%%%%%%%%%%%%%%%%%%%%%%%%%%%%%%

\paragraph{Relation to Dual-Process Theory.}
The Awareness framework extends Kahneman's System 1/System 2:

\begin{table}[htbp]
\centering
\caption{Awareness vs. Dual-Process Theory}
\label{tab:dual-process-comparison}
\renewcommand{\arraystretch}{1.3}
\begin{tabular}{@{}lll@{}}
\toprule
\textbf{Kahneman} & \textbf{Awareness} & \textbf{Extension} \\
\midrule
System 1 & SYS 0 + SYS 1 & Distinguishes visceral from emotional \\
System 2 & SYS 2 & Same: reflective, effortful \\
--- & Self vs. Community & Adds social dimension \\
--- & INU/KNU/IDN & Adds utility category specificity \\
--- & 80 Axioms & Provides micro-mechanisms \\
\bottomrule
\end{tabular}
\end{table}

\paragraph{Relation to Prospect Theory.}
Key Prospect Theory findings emerge from awareness asymmetries:

\begin{itemize}
\item \textbf{Loss aversion:} $A_{SYS1}(fear) > A_{SYS1}(hope)$ (AWX-A16 vs. A17)
\item \textbf{Reference dependence:} $A$ is calibrated to reference point
\item \textbf{Diminishing sensitivity:} $A$ saturates at extremes
\item \textbf{Probability weighting:} $A \propto Availability$, not objective probability
\end{itemize}

\paragraph{Relation to Construal Level Theory.}
Psychological distance effects map to awareness:

\begin{itemize}
\item \textbf{Temporal distance:} $A_{SYS0}$ (near) vs. $A_{SYS2}$ (far)
\item \textbf{Social distance:} $A^{self}$ (near) vs. $A^{comm}$ (far)
\item \textbf{Hypothetical distance:} $A_{SYS0}$ (real) vs. $A_{SYS2}$ (hypothetical)
\end{itemize}

%%%%%%%%%%%%%%%%%%%%%%%%%%%%%%%%%%%%%%%%%%%%%%%%%%%%%%%%%%%%%%%%%%%%%%%%%%%%%%%
\subsubsection{Summary: Awareness as Bias Unifier}
%%%%%%%%%%%%%%%%%%%%%%%%%%%%%%%%%%%%%%%%%%%%%%%%%%%%%%%%%%%%%%%%%%%%%%%%%%%%%%%

\begin{tcolorbox}[colback=gray!5!white, colframe=gray!75!black, title=Key Insight: Biases as Awareness Asymmetries]

\textbf{The Awareness framework provides a unified explanation for 59+ documented biases:}

\textbf{Six Sources of Asymmetry:}
\begin{enumerate}
\item \textbf{System asymmetry:} SYS 0/1 faster than SYS 2
\item \textbf{Self-community asymmetry:} $A^{self} >> A^{comm}$ (gradient)
\item \textbf{Category asymmetry:} INU instant, KNU never instant
\item \textbf{Determinant asymmetry:} Available > Base rate; Fear > Hope
\item \textbf{Temporal asymmetry:} Decay, Habituation
\item \textbf{Defensive asymmetry:} Ego-consistent > Ego-threatening
\end{enumerate}

\textbf{The Debiasing Implication:}
$$\text{Debias} = \text{Activate underactivated } A \text{ / Dampen overactivated } A$$

\textbf{The Theoretical Contribution:}
Awareness provides a \emph{generative} framework that:
\begin{itemize}
\item Explains existing biases
\item Predicts new biases
\item Guides debiasing interventions
\item Integrates with dual-process, prospect, and construal theories
\end{itemize}

\end{tcolorbox}

%%%%%%%%%%%%%%%%%%%%%%%%%%%%%%%%%%%%%%%%%%%%%%%%%%%%%%%%%%%%%%%%%%%%%%%%%%%%%%%
% END OF BIAS SECTION
%%%%%%%%%%%%%%%%%%%%%%%%%%%%%%%%%%%%%%%%%%%%%%%%%%%%%%%%%%%%%%%%%%%%%%%%%%%%%%%
